% Section 6, part 4.
\begin{lem}[Functions of $x$ pair to zero]
\label{lem:mod-xfun}
Let $2h-2+n>0$, $n\geq 1$, let $p_2,\dots,p_n$ be points of $\bbP^1$ that are not branch points, and let $f$ be holomorphic near the five critical values $x(\varsigma_b)$ (for instance $f\in\bbQ(v)[X]$, $v\neq 0$ fixed). Then for each $b$,
\[
\Res_{p=\varsigma_b}f(x(p))\,\omega_{h,n}(p,p_2,\dots,p_n)=0 .
\]
Consequently every pairing of $\omega_{h,n}$ in which one slot is paired with $f(x)$ by the residue pairing, the other slots being paired by residues at the branch points, on circles or along rays, vanishes; legs $\phi,\phi'\in\cL$ with $\phi-\phi'\in\bbQ(v)[x]$ have the same residue pairings; and (with $f=1$) $\omega_{h,n}$ has zero residues at the branch points in each variable.
\end{lem}

\begin{proof}
\emph{The trace identity.} For $X\in\bbP^1$ let $p^1(X),\dots,p^5(X)$ be the roots of $T^5-4XT+4v$, the points of $x^{-1}(X)$. We claim $\mathrm{Tr}\,\omega(X):=\sum_{i=1}^5\omega_{h,n}(p^i(X),p_2,\dots,p_n)=0$. For $q'\notin x^{-1}(X)$, $\prod_i(q'-p^i(X))=4q'\,(x(q')-X)$, so $\sum_i\omega_{0,2}(p^i(X),q')=\dd X\,\dd x(q')/(X-x(q'))^2$, and hence $\sum_i\int_{\iota_bq}^q\omega_{0,2}(p^i(X),\cdot)=\dd X\big[(X-x(q))^{-1}-(X-x(\iota_bq))^{-1}\big]=0$. Take disjoint $\iota_b$-stable small discs $U_b\ni\varsigma_b$ avoiding $p_2,\dots,p_n$, $U:=\bigcup_bU_b$. For $p_1\notin U$, \eqref{eq:mod-kernel} writes $\omega_{h,n}(p_1,J)$ as $\sum_b(2\pi\ri)^{-1}\oint_{\partial U_b}K_b(p_1,q)$ times the bracket, which does not depend on $p_1$ (the path in $K_b$ stays in $U_b$). For $X\notin x(\overline U)$, summing over $i$ under the integrals gives $\mathrm{Tr}\,\omega(X)=0$, hence $\mathrm{Tr}\,\omega=0$ by meromorphy (as in \cite[App.~A, (A-1), (A-5)]{Eynard:2007kz}).

\emph{The residues.} Put $X_b:=x(\varsigma_b)$. The fibre $x^{-1}(X_b)$ is $\varsigma_b$ (local degree $2$) and three points $q_1,q_2,q_3$ of local degree $1$, none a branch point (the critical values are distinct) and none $0$ or $\infty$. So $\eta:=f(x)\,\omega_{h,n}(\cdot,J)$ is holomorphic at the $q_k$ (\cref{conv:mod-kernel}). The preimage of a small circle around $X_b$ is a union of loops around the points of the fibre, so $\Res_{X_b}\mathrm{Tr}(\eta)=\sum_{q\in x^{-1}(X_b)}\Res_q\eta$, and $\mathrm{Tr}(\eta)=f(X)\,\mathrm{Tr}\,\omega=0$ gives $\Res_{\varsigma_b}\eta=0$.

\emph{Consequences.} By symmetry the slot carrying $f(x)$ may be the first; the contours of the other slots can be taken disjoint from small discs around the branch points, and there are no diagonal poles, so the residue in the first slot may be taken first, and it vanishes. (Cf.\ \cite[Thm.~4.4 (4-23), Cor.~4.1 (4-26)--(4-29), App.~A (A-25)]{Eynard:2007kz}, not used here.)
\end{proof}

\subsection{The leg dictionary and the residue identity}
\label{sec:mod-legs}

\begin{defn}[The legs]
\label{def:mod-legs}
Let $\phi^b_0:=p,\ \log p,\ v/p,\ p^3,\ p^2$ for $b=0,\dots,4$. For $F$ in the span of the $p^j(\log p)^l$ ($j\in\bbZ$, $l\geq 0$) with coefficients in $\bbQ[v]$, let $I(F)$ be the termwise antiderivative of $-F\,\dd x=-F(p^3-vp^{-2})\,\dd p$ without a term $p^0(\log p)^0$ (that of $p^j(\log p)^l\dd p$ is $p^{j+1}$ times a polynomial in $\log p$ if $j\neq -1$, and $(\log p)^{l+1}/(l+1)$ if $j=-1$). Put $\phi^b_{k+1}:=I(\phi^b_k)$. Then $D_x\phi^b_{k+1}=-\phi^b_k$, and $\phi^b_k$ is homogeneous of weight $\wt(\phi^b_k)=w_b+4k$, where $w_b:=\nu_b+1$, that is $(w_0,\dots,w_4)=(1,0,4,3,2)$.
\end{defn}

\begin{prop}[The leg dictionary]
\label{prop:mod-LD}
Let $\varrho=(\varrho_0,\dots,\varrho_4):=(-1,1,-1,-3,-2)$. For all $a,b\in\bbZ/5$ and $k,d\geq 0$,
\begin{equation}
\label{eq:mod-LD}
\langle\phi^b_k,\dd\xi_{a,d}\rangle=\sum_c\Res_{c}\phi^b_k\,\dd\xi_{a,d}=-\varrho_b\,\delta_{ab}\,\delta_{dk},
\end{equation}
$c$ running over the branch points and $\log p$ being any branch near each of them. Hence $q_{\phi^b_k}=-\varrho_b\,z^kH^b$, and the TR leg of $v^mz^kH^b$ (the insertion $v^m\psi^kH^b$, \cref{conv:tra-forms}) is $-\varrho_b^{-1}v^m\phi^b_k$. In particular the unit $H^0$ has the TR leg $y=p$; $\Phi_{y\dd x}=p^5/5-v\log p=-\phi^0_1$ has $q_{\Phi_{y\dd x}}=-z\,H^0$; $-\log p$ has $q_{-\log p}=H$; the Euler-dilaton vector $\se=z\,H^0+vH$ has the TR leg $-p^5/5$; and $q_{f(x)}=0$ for $f\in\bbQ(v)[x]$.
\end{prop}

\begin{proof}
\emph{Step 1 (branches and functions of $x$).} Each leg is analytic near every $c\neq 0$. A change of branch $\log p\mapsto\log p+2\pi\ri N$ changes $\phi^b_k$ by a polynomial $\delta^b_k$ in $x$ with constant coefficients, by induction: $\delta^b_0\in\{0,2\pi\ri N\}$, and $\dd\phi^b_{k+1}=-\phi^b_k\,\dd x$ survives the substitution, so $\delta^b_{k+1}+\int_0^x\delta^b_k$ is constant. For $f$ holomorphic near $c$, repeated integration by parts gives $\Res_cf(x)\,\dd\xi_{a,d}=\pm\Res_cf^{(d+1)}(x)\,\xi_a\dd x=0$, as $\xi_a\dd x=\varepsilon'_av^{\alpha_a}p^{m_a-2}\dd p$ is holomorphic at $c$. So the pairings do not depend on the branches, and $q_{f(x)}=0$.
\emph{Step 2 (lowering).} For $d\geq 1$, $\xi_{a,d}\dd\phi=(D_x\phi)\,\dd\xi_{a,d-1}$, so integration by parts gives $\langle\phi,\dd\xi_{a,d}\rangle=-\langle D_x\phi,\dd\xi_{a,d-1}\rangle$, that is $\langle\phi^b_{k+1},\dd\xi_{a,d}\rangle=\langle\phi^b_k,\dd\xi_{a,d-1}\rangle$.
\emph{Step 3 ($d=0$, $k\geq 1$).} $\langle\phi^b_k,\dd\xi_a\rangle=\sum_c\Res_c\phi^b_{k-1}\,\xi_a\dd x=0$, as $\xi_a\dd x$ is holomorphic at $c$.
\emph{Step 4 ($k=0$, $d\geq 1$).} $\langle\phi^b_0,\dd\xi_{a,d}\rangle=-\sum_c\Res_cD_x^d(\xi_a)\,\dd\phi^b_0$. This rational form is holomorphic at $p=0$ ($D_x$ maps functions holomorphic at $0$ to $O(p^2)$, and $\dd\phi^b_0=O(p^{-2})\dd p$) and $O(p^{1-4d})\dd p$ at $\infty$; by the residue theorem the sum vanishes.
\emph{Step 5 ($k=d=0$).} Write $\dd\phi^b_0=\gamma_bp^{n_b}\dd p$ with $(\gamma_b,n_b)=(1,0),(1,-1),(-v,-2),(3,2),(2,1)$ (local constants). Since $\Res_cp^N\dd p/(p^5-v)=c^{N-4}/5$ and $\sum_cc^N=5v^{N/5}$ if $5\mid N$ and $0$ otherwise, $\langle\phi^b_0,\dd\xi_a\rangle=-\varepsilon'_a\gamma_bv^{\alpha_a}\,\frac15\sum_cc^{m_a+n_b-4}$, non-zero only if $m_a+n_b\equiv 4\pmod 5$, that is $a=b$; for $a=b=0,\dots,4$ the values are $1,-1,1,3,2=-\varrho_b$.
\emph{Step 6.} If $k=d$, Steps 2 and 5 give $-\varrho_b\delta_{ab}$; if $k>d$, Steps 2 and 3 give $0$; if $k<d$, Steps 2 and 4 give $0$. This is \eqref{eq:mod-LD}. The consequences: $\phi^0_0=p$, $-\varrho_0=1$; $\phi^0_1=-(p^5/5-v\log p)$, so $q_{\Phi_{y\dd x}}=\varrho_0z\,H^0=-zH^0$; $q_{\log p}=-\varrho_1H=-H$; and the TR leg of $\se$ is $\phi^0_1-v\phi^1_0=-p^5/5$.
\end{proof}

\begin{cor}[Free energies of the model and of the curve]
\label{cor:mod-trF}
For $h\geq 2$ let $\Fm_h(v)=\langle\ \rangle_h(v)=\int_{\overline{\mathcal M}_h}\Omega^{\mathrm{mod}}_{v;h,0}$ be the genus-$h$ free energy of the model CohFT (\cref{def:set-modcohft}), the value at the base point $-z1$ of the dilaton-shifted potential $\overline F_h$ (\cref{def:set-modcohft,def:set-graphsum}). Then $\Fm_h=(-5)^{1-h}\FTR_h$.
\end{cor}

\begin{proof}
By \cref{conv:mod-kernel}, $(2-2h)\FTR_h=\langle\Phi_{y\dd x}\rangle^{\mathrm{TR}}_h$, for any branch of $\log p$ (\cref{lem:mod-xfun}); only this definition is used (cf.\ \cite[Def.~4.3 (4-10), Thm.~4.7 (4-30)]{Eynard:2007kz}). By \cref{thm:mod-tr} and \cref{prop:mod-LD}, $\langle\Phi_{y\dd x}\rangle^{\mathrm{TR}}_h=-(-5)^{h-1}\langle\psi H^0\rangle_h$, and the dilaton equation for the flat unit $H^0$ (\cref{lem:set-dilaton}~(i), (ii)) gives $\langle\psi H^0\rangle_h=(2h-2)\Fm_h$. Comparing, $(2-2h)\FTR_h=(2-2h)(-5)^{h-1}\Fm_h$.
\end{proof}

% Section 6, part 5.
\begin{defn}[Weight spaces, the leg map, the leg class]
\label{def:mod-weights}
For $w\geq 0$ let $V_w$ be the $\bbQ$-span of the insertions $v^m\psi^kH^b$ ($m,k\geq 0$, $b\in\bbZ/5$) with $w_b+4k+5m=w$, and $\mathbb F:=\bigoplus_wV_w$, the free graded $\bbQ[v]$-module on the $\psi^kH^b$ (the weights of \cref{conv:tra-forms}). The \emph{leg map} is the $\bbQ$-linear map $\Lambda\colon\mathbb F\to\{\text{legs}\}$, $\Lambda(v^m\psi^kH^b):=\varrho_b^{-1}v^m\phi^b_k$; by \cref{prop:mod-LD} the TR leg of $X\in\mathbb F$ is $-\Lambda(X)$. Let $\bbQ[p,\bar p]$, $\bar p=v/p$, be the subring of $\bbQ[v][p,p^{-1}]$ spanned by the $v^mp^j$ with $j\geq -m$; recall the leg class $\cL=\bbQ[v][p]+\bbQ(v)[x]$ (\cref{def:mod-pair}). Put
\[
\mathsf M:=\bbQ[p,\bar p]\oplus\log p\cdot\bbQ[v,x],\qquad \mathsf N:=\mathsf M/\bbQ[v,x]=\bbQ[p,\bar p]/\bbQ[v,x]\oplus\log p\cdot\bbQ[v,x],
\]
(the sum is direct since $1$ and $\log p$ are independent over Laurent polynomials; $\bbQ[v,x]\subset\bbQ[p,\bar p]$ as $v=p\bar p$, $x=p^4/4+\bar p$), graded by weight. For $F=A+\log p\cdot B\in\mathsf M$ the element $B\in\bbQ[v,x]$ is the \emph{$\log p$-coefficient} of $F$, and $F$ is \emph{log-free} if $B=0$. Finally $\mathrm d(w):=\lfloor w/5\rfloor+1-[5\mid w]$ and $\mathrm e(w):=\#\{(m,n)\in\bbZ_{\geq 0}^2: 5m+4n=w\}$.
\end{defn}

\begin{lem}[Closure]
\label{lem:mod-closure}
\begin{enumerate}
\item For $a,c\geq 0$, $I(p^a\bar p^c)=-p^{a+4}\bar p^c/(a-c+4)+p^a\bar p^{c+1}/(a-c-1)$, a term with a zero denominator being read as $-v^c\log p$, respectively $+v^{c+1}\log p$.
\item For $g\in\bbQ[v,x]$ with an $x$-primitive $G$, $I(\log p\cdot g(x))=-\log p\cdot G(x)+\int G(x)\,\dd p/p$, where the last integral is termwise: $\int p^a\bar p^c\,\dd p/p=p^a\bar p^c/(a-c)$ if $a\neq c$ and $=v^c\log p$ if $a=c$.
\item $I$ is $\bbQ[v]$-linear, $I(\mathsf M)\subset\mathsf M$ and $I(\bbQ[v,x])\subset\bbQ[v,x]$. Hence every leg $\phi^b_k$ lies in $\mathsf M$, and $\mathsf M$ is a $\bbQ[v,x]$-module.
\end{enumerate}
\end{lem}

\begin{proof}
(i) $-\int v^cp^{a-c}(p^3-vp^{-2})\,\dd p=-v^cp^{a-c+4}/(a-c+4)+v^{c+1}p^{a-c-1}/(a-c-1)$, with $\int\dd p/p=\log p$ in the exceptional cases. (ii) $\dd(\log p\cdot G(x))=G(x)\,\dd p/p+\log p\cdot g(x)\,\dd x$ with $G(x(p))\in\bbQ[p,\bar p]$, and the right side has no term $p^0(\log p)^0$, so it is the antiderivative of \cref{def:mod-legs}; replacing $G$ by $G+C$ changes the two terms by $\mp C\log p$. (iii) By (i), $I(\bbQ[p,\bar p])\subset\mathsf M$; by (ii), $I(\log p\cdot g)\in\mathsf M$, as $\int G\,\dd p/p\in\bbQ[p,\bar p]+\log p\cdot\bbQ[v]$. For $g\in\bbQ[v,x]$, $-g(x)\,\dd x=-\dd G(x(p))$ has no term $p^{-1}\dd p$, so $I(g)=-G(x)+[p^0]G(x(p))\in\bbQ[v,x]$. The $\phi^b_0$ lie in $\mathsf M$, so all $\phi^b_k$ do by induction; and $\bbQ[p,\bar p]$ is a ring containing $\bbQ[v,x]$.
\end{proof}

\begin{lem}[Generating function of the log coefficients]
\label{lem:mod-genfn}
For all $b$ and $k$, the $\log p$-coefficient of $\phi^b_k$ is $\sum_{i=0}^k\mathcal R^b_{k-i}(-x)^i/i!$, where $\mathcal R^b_n:=\Res_{p=0}\Wv^n\,\dd\phi^b_0/n!$, a monomial in $v$.
\end{lem}

\begin{proof}
Fix $b$, write $\phi_k:=\phi^b_k$, and put $\mathcal B_n:=\sum_{j=0}^nx^j\phi_{n-j}/j!$, the coefficient of $\epsilon^n$ in $e^{\epsilon x}\sum_k\epsilon^k\phi_k$. Since $\dd\phi_m/\dd p=-x'\phi_{m-1}$ for $m\geq 1$, $\dd\mathcal B_n/\dd p=x^n\phi_0'/n!$ telescopes to a Laurent polynomial in $p$. By \cref{lem:mod-closure}~(iii), $\mathcal B_n=A+\log p\cdot B$ with $A\in\bbQ[p,\bar p]$, $B\in\bbQ[v,x]$; then $\dd\mathcal B_n/\dd p=A'+B/p+\log p\cdot\dd B(x(p))/\dd p$ has no logarithmic term, so $B\in\bbQ[v]$, and as $A'$ has no term $p^{-1}$, $B=[p^{-1}]x^n\phi_0'/n!=\mathcal R^b_n$. Inverting, $\phi_k=\sum_{i=0}^k(-x)^i\mathcal B_{k-i}/i!$, and taking $\log p$-coefficients, which commutes with multiplication by $\bbQ[v,x]$, gives the claim. By weight, $\mathcal R^b_n$ is a monomial in $v$.
\end{proof}

\begin{lem}[The residue identity]
\label{lem:mod-resid}
For all $n\geq 0$ and $b$, $\mathcal R^b_n=\varrho_b\,\Pi_0(\psi^nH^b)$.
\end{lem}

\begin{proof}
By \cref{def:gau-van}, $\IC_m=\Res_{p=0}p^me^{-\Wv(p)/z}\dd p$ coefficientwise; with $\YC_b=\varepsilon_b\IC_{\nu_b}$ and $\Pi_0(z^nH^b)=(-1)^{n+1}[z^{-n}]\YC_b$ (\cref{def:tra-tower}) this gives
\[
[z^{-n}]\IC_m=\frac{(-1)^n}{n!}\Res_{p=0}p^m\Wv^n\,\dd p,\qquad \Pi_0(\psi^nH^b)=-\frac{\varepsilon_b}{n!}\Res_{p=0}p^{\nu_b}\Wv^n\,\dd p .
\]
Compare with $\dd\phi^b_0=\dd p,\ \dd p/p,\ -vp^{-2}\dd p,\ 3p^2\dd p,\ 2p\,\dd p$: for $b=0,1,3,4$ the identity is immediate, with $\varrho_b=-1,1,-3,-2$. For $b=2$, $\Res_{p=0}\Wv^n(p^3-vp^{-2})\,\dd p=\Res_{p=0}\dd(\Wv^{n+1})/(n+1)=0$, so $\mathcal R^2_n=-\Res_{p=0}p^3\Wv^n\dd p/n!=-\Pi_0(\psi^nH^2)$, and $\varrho_2=-1$.
\end{proof}

\begin{prop}[Log coefficients of the legs]
\label{prop:mod-dii}
For all $b$ and $k$,
\begin{equation}
\label{eq:mod-dii}
\text{the }\log p\text{-coefficient of }\phi^b_k\ =\ \varrho_b\sum_{i=0}^k\Pi_i(\psi^kH^b)\,\frac{(-x)^i}{i!}.
\end{equation}
Hence, for $X\in V_w$, the $\log p$-coefficient of $\Lambda(X)$ is $\sum_{i\geq 0}\Pi_i(X)(-x)^i/i!$, and $\Lambda(X)$ is log-free if and only if $X\in\sKo$.
\end{prop}

\begin{proof}
\Cref{lem:mod-genfn,lem:mod-resid} and the shift rule $\Pi_i(\psi^kH^b)=\Pi_0(\psi^{k-i}H^b)$ of \cref{def:tra-tower} give \eqref{eq:mod-dii}. For $X\in V_w$ the factors $\varrho_b^{-1}$ of $\Lambda$ cancel, and the $\Pi_i$ are $\bbQ[v]$-linear. As $\Pi_i(X)\in\bbQ[v]$ and the powers of $x$ are independent over $\bbQ(v)$, $\Lambda(X)$ is log-free if and only if $\Pi_i(X)=0$ for all $i$, that is $X\in\sKo$ (\cref{def:tra-slice}).
\end{proof}

% Section 6, part 6.
\subsection{The lattice theorem}
\label{sec:mod-lattice}

\begin{lem}[Counting]
\label{lem:mod-count}
For every $w\geq 0$, $V_w$ has exactly $\mathrm d(w)$ basis triples $(m,k,b)$ with $b\neq 1$ and exactly $\mathrm e(w)$ with $b=1$.
\end{lem}

\begin{proof}
Fix $m$ with $5m\leq w$ and put $r:=w-5m$. A triple with $b\neq 1$ needs $w_b+4k=r$, $w_b\in\{1,4,3,2\}$ for $b=0,2,3,4$; these meet every class modulo $4$ once, so there is one candidate $b$, which fails only if $w_b>r$, that is $r=0$, $w=5m$. This gives $\mathrm d(w)$. The triples with $b=1$ ($w_1=0$) are the pairs $(m,k)$ with $5m+4k=w$.
\end{proof}

\begin{lem}[The leg class]
\label{lem:mod-classL}
\begin{enumerate}
\item $\bbQ[p,\bar p]=\bbQ[p,x]=\bbQ[v][p]+\bbQ[v][x]$, and $\bbQ[p,x]$ is a free $\bbQ[v,x]$-module with basis $1,p,p^2,p^3,p^4$.
\item $\bbQ[v][p^{\pm1}]\cap\bbQ(v)[x]=\bbQ[v][x]$.
\item $\cL\cap\bbQ[v][p^{\pm1}]=\bbQ[p,\bar p]$: a Laurent polynomial with coefficients in $\bbQ[v]$ lies in $\cL$ if and only if each of its terms $v^mp^j$ has $j\geq -m$.
\item $\cL/\bbQ(v)[x]=\bbQ[p,\bar p]/\bbQ[v,x]=\bbQ[v][p]/\bbQ[v]=\bigoplus_{i=1}^4\bbQ[v,x]\,p^i$; it is free over $\bbQ[v]$, of dimension $\mathrm d(w)$ in weight $w$.
\end{enumerate}
\end{lem}

\begin{proof}
(i) $\bar p=x-p^4/4$ gives $\bbQ[p,\bar p]=\bbQ[p,x]$, and $v=px-p^5/4$, so $p,x$ are algebraically independent. The map $\bbQ[v,x][T]/(T^5-4xT+4v)\to\bbQ[p,x]$, $T\mapsto p$, is onto, and its source is $\bbQ[x,T]$, so it is an isomorphism, which gives the basis. Clearly $\bbQ[v][p]+\bbQ[v][x]\subset\bbQ[p,\bar p]$. Conversely, if $F\in\bbQ[p,\bar p]$ is homogeneous and $v^mp^{-n}$ ($n>0$, $m\geq n$) is its term with the most negative exponent, unique by homogeneity, subtract a multiple of $v^{m-n}x^n$, whose lowest term is $v^np^{-n}$, and induct. (ii) If $P=\sum_{n\leq N}a_nx^n\in\bbQ[v][p^{\pm1}]$ with $a_N\neq 0$, only $x^N$ contributes to $p^{4N}$, with coefficient $a_N4^{-N}$, so $a_N\in\bbQ[v]$; induct on $N$. (iii) $\supset$ follows from (i). If $F=R+P$, $R\in\bbQ[v][p]$, $P\in\bbQ(v)[x]$, is a Laurent polynomial over $\bbQ[v]$, then $P\in\bbQ[v][x]$ by (ii), and $F\in\bbQ[p,\bar p]$ by (i). (iv) $\bbQ[p,\bar p]\to\cL/\bbQ(v)[x]$ is onto, with kernel $\bbQ[p,\bar p]\cap\bbQ(v)[x]=\bbQ[v,x]$ by (ii); the quotient is $\bigoplus_{i=1}^4\bbQ[v,x]p^i$ by (i). As the lowest term of $x^n$ is $v^np^{-n}$, $\bbQ[v][p]\cap\bbQ[v][x]=\bbQ[v]$, so the quotient is also $\bbQ[v][p]/\bbQ[v]$, with the basis $v^mp^{w-5m}$, $5m<w$, in weight $w$: $\mathrm d(w)$ elements.
\end{proof}

\begin{lem}[Reduction modulo $v$]
\label{lem:mod-modv}
In every weight, $\Lambda\otimes\bbQ[v]/(v)\colon\mathbb F/v\mathbb F\to\mathsf N/v\mathsf N$ is an isomorphism; here $\Lambda$ is followed by the projection $\mathsf M\to\mathsf N$.
\end{lem}

\begin{proof}
$\Lambda$ is graded, $\bbQ[v]$-linear, with values in $\mathsf M$ (\cref{lem:mod-closure}). \emph{Source.} $(\mathbb F/v\mathbb F)_w$ has the basis $\psi^kH^b$ with $w_b+4k=w$: one element if $4\nmid w$; $\{\psi^{n-1}H^2,\psi^nH\}$ if $w=4n\geq 4$; $\{H\}$ if $w=0$. \emph{Target.} $\mathsf N/v\mathsf N=\bbQ[p,\bar p]/(\bbQ[v,x]+p\bar p\,\bbQ[p,\bar p])\oplus\log p\cdot\bbQ[x]$. Here $\bbQ[p,\bar p]/(p\bar p)$ has the basis $1$, $p^a$, $\bar p^c$ ($a,c\geq 1$), and $\bbQ[v,x]$ maps to the span of $1$ and $x^n\equiv p^{4n}/4^n+\bar p^n$; so $(\mathsf N/v\mathsf N)_w$ has the basis $\{p^w\}$ if $w\geq 1$, $4\nmid w$; $\{p^{4n},\log p\cdot x^n\}$ if $w=4n\geq 4$; $\{\log p\}$ if $w=0$. \emph{The map.} By \cref{lem:mod-closure}~(i), $I(p^a)\equiv -p^{a+4}/(a+4)$ mod $v\mathsf M$ ($a\geq 1$), so $\phi^b_k\equiv(-1)^kp^{w_b+4k}/\prod_{i=1}^k(w_b+4i)$ for $b=0,4,3$; as $\bar p\equiv -p^4/4$ modulo $\bbQ[v,x]$, $\phi^2_k\equiv(-1)^{k+1}p^{4k+4}/\big(4\prod_{i=1}^k(4+4i)\big)$; and by \cref{lem:mod-closure}~(ii) the $\log p$-component of $\phi^1_k$ is $(-1)^kx^k/k!$ modulo $v$. In weight $4n$ the matrix is triangular with non-zero diagonal, in the other weights a non-zero $1\times 1$ matrix.
\end{proof}

\begin{thm}[The lattice theorem]
\label{thm:mod-lattice}
For every $w\geq 0$:
\begin{enumerate}
\item[(a)] $V_w$ has $\mathrm d(w)$ basis triples with $b\neq 1$ and $\mathrm e(w)$ with $b=1$;
\item[(b)] for $X\in\sKo\cap V_w$, $\Lambda(X)$ lies in $\bbQ[p,\bar p]=\bbQ[v][p]+\bbQ[v][x]$, exactly and not only modulo functions of $x$;
\item[(c)] $\Lambda\colon\sKo\cap V_w\to(\cL/\bbQ(v)[x])_w$ is bijective, and both sides have dimension $\mathrm d(w)$;
\item[(d)] $\Lambda\colon V_w\to\mathsf N_w$ is bijective; $\mathsf N$ is a free $\bbQ[v]$-module with basis $\{\Lambda(\psi^kH^b)\}$.
\end{enumerate}
The coefficients must lie in $\bbQ[v]$: over $\bbQ(v)$ the statement fails, since $\Lambda(H^2/v)=-1/p\notin\cL$.
\end{thm}

\begin{proof}
(a) is \cref{lem:mod-count}. (b) $\Lambda(X)\in\mathsf M$ (\cref{lem:mod-closure}) is log-free (\cref{prop:mod-dii}), so it lies in $\bbQ[p,\bar p]$, and \cref{lem:mod-classL}~(i) applies. (d) Let $C$ be the cokernel of the graded $\bbQ[v]$-linear map $\Lambda\colon\mathbb F\to\mathsf N$, whose pieces are finite-dimensional and zero in negative weights. By right exactness and \cref{lem:mod-modv}, $C/vC=0$; if $w$ were minimal with $C_w\neq 0$, then $(vC)_w=vC_{w-5}=0$ and $C_w=0$. So $\Lambda$ is onto, and as $\dim V_w=\mathrm d(w)+\mathrm e(w)=\dim\mathsf N_w$ (\cref{lem:mod-count}, \cref{lem:mod-classL}~(iv), and $\dim\bbQ[v,x]_w=\mathrm e(w)$), $\Lambda_w$ is bijective. (c) By (b) and \cref{lem:mod-classL}~(iv) the map is well defined. If $\Lambda(X)\in\bbQ(v)[x]$, then $\Lambda(X)\in\bbQ[v,x]$ (\cref{lem:mod-classL}~(ii)), so $\Lambda(X)=0$ in $\mathsf N$ and $X=0$ by (d). Given $\phi\in(\bbQ[p,\bar p]/\bbQ[v,x])_w$, $X:=\Lambda^{-1}(\phi\oplus 0)$ has a log-free leg, so $X\in\sKo$ (\cref{prop:mod-dii}). The dimensions are \cref{lem:mod-classL}~(iv). Finally $\Lambda(H^2)=\varrho_2^{-1}v/p=-v/p$, and $-1/p\notin\cL$ by \cref{lem:mod-classL}~(iii).
\end{proof}

\begin{cor}[Splitting of $V_w$]
\label{cor:mod-split}
$V_w=(\sKo\cap V_w)\oplus\operatorname{span}\{v^m\psi^kH:5m+4k=w\}$. Every $X\in V_w$ is $X_0+\sum a_{m,k}v^m\psi^kH$ with $X_0\in\sKo$, the $a_{m,k}$ being read off from the $\log p$-coefficient of $\Lambda(X)$.
\end{cor}

\begin{proof}
By \cref{lem:mod-closure}~(ii) and induction, the $\log p$-coefficient of $\phi^1_k$ is $(-1)^kx^k/k!$ plus terms of lower degree in $x$, so $X\mapsto$ the $\log p$-coefficient of $\Lambda(X)$ is triangular, hence bijective, from the span of the $v^m\psi^kH$ onto $\bbQ[v,x]_w$ ($v^m\psi^kH\leftrightarrow v^mx^k$); its kernel on $V_w$ is $\sKo\cap V_w$ by \cref{prop:mod-dii}.
\end{proof}

\begin{cor}[Splitting of $\sK\cap V_w$]
\label{cor:mod-Ksplit}
$\sK\cap V_w=(\sKo\cap V_w)\oplus\bbQ v^{w/5}H$ if $5\mid w$, and $\sK\cap V_w=\sKo\cap V_w$ otherwise. On $\sK$ the only legs beyond $\bbQ[p,\bar p]$ are multiples of $\log p$: for $X\in\sK\cap V_w$, $\Lambda(X)=\Lambda(X-\Pi_0(X)H)+\Pi_0(X)\log p$ with $\Lambda(X-\Pi_0(X)H)\in\bbQ[p,\bar p]$.
\end{cor}

\begin{proof}
By \cref{prop:mod-dii}, $X\in\sK$ if and only if the $\log p$-coefficient of $\Lambda(X)$ equals $\Pi_0(X)\in\bbQ v^{w/5}$ (\cref{def:tra-tower}; zero unless $5\mid w$). Write $X=X_0+Y$ as in \cref{cor:mod-split}; $\Lambda(Y)$ and $\Lambda(\Pi_0(X)H)=\Pi_0(X)\log p$ have the same $\log p$-coefficient, so $Y=\Pi_0(X)H$ by injectivity on the span of the $v^m\psi^kH$. Conversely $\sKo\subset\sK$ and $H\in\sK$. The last statement is \cref{thm:mod-lattice}~(b) for $X-\Pi_0(X)H\in\sKo\cap V_w$.
\end{proof}

% Section 6, part 7.
\subsection{L-regularity and the Gap Theorem}
\label{sec:mod-Lreg}

We use the closed formula of \cite[Thm.~7.1]{ABDKS25}: for a rational function $x(z)$ on $\bbP^1$ such that $\dd x$ has only simple zeros, the differentials of topological recursion for $(x=x(z),y=z)$ with the Bergman kernel and $\omega^{(0)}_{1,0}=-y\,\dd x$ are
\begin{equation}
\label{eq:mod-abdks}
\omega^{(g)}_{m,0}=[\hbar^{2g}]\sum_{r_1,\dots,r_m\geq 0}\prod_{i=1}^m\big(\mathcal D_{x_i}\big)^{r_i}[w_i^{r_i}]\sum_{\Gamma}\hbar^{2g(\Gamma)}\prod_{\{i,j\}\in E(\Gamma)}W_{i,j}\prod_{i=1}^mW_i
\end{equation}
\cite[(7.9)]{ABDKS25}, where $\Gamma$ runs over the connected simple graphs on the vertices $1,\dots,m$, $g(\Gamma)$ is the first Betti number, $\mathcal D_{x_i}\eta:=\dd(\eta/\dd x_i)$ on $1$-forms in $z_i$, $S(u):=\sinh(u/2)/(u/2)$,
\begin{gather*}
W_i:=-\exp\!\big(-w_i(S(w_i\hbar\partial_{z_i})-1)\,x(z_i)\big)\frac{\dd z_i}{w_i},\\
W_{i,j}:=\frac{1}{\hbar^2}\Big(\exp\!\big(\hbar^2w_iw_jS(w_i\hbar\partial_{z_i})S(w_j\hbar\partial_{z_j})(z_i-z_j)^{-2}\big)-1\Big)
\end{gather*}
\cite[(7.6), (7.7)]{ABDKS25}, and $W_{i,j}=\sum_{k\geq 0}(\hbar^2/4)^k(w_i+w_j)^{2k}w_iw_j(z_i-z_j)^{-2k-2}$; for one point \cite[(7.10)]{ABDKS25},
\begin{equation}
\label{eq:mod-710}
\sum_{g\geq 1}\hbar^{2g}\omega^{(g)}_{1,0}=\sum_{r\geq 0}\mathcal D_x^r[w^r]W_1 .
\end{equation}
The standing assumptions of \cite[\S1.1]{ABDKS25} (simple critical points of $x$, $\dd y\neq 0$ there) hold for $x=\Wv$, $y=p$, $v\neq 0$ (\cref{def:mod-curve}), and $\omega_{h,n}=(-1)^n\omega^{(h)}_{n,0}$ (\cref{conv:mod-kernel}).

\begin{lem}[Structure of the correlators]
\label{lem:mod-abdks}
For $2h-2+n>0$ there are an integer $K\geq 0$ and $N_{h,n}\in\bbQ[v][p_1,\dots,p_n]$ with
\[
\omega_{h,n}=N_{h,n}(p;v)\,\frac{\prod_i\dd p_i}{\prod_i(p_i^5-v)^K}.
\]
\end{lem}

\begin{proof}
(1) \emph{For fixed $(h,n)$, \eqref{eq:mod-abdks} is a finite sum of terms in $\bbQ[v][p_i,p_i^{-1},(p_i^5-v)^{-1},(p_i-p_j)^{-1}]\prod_i\dd p_i$.} Since $S(u)-1=\sum_{k\geq 1}u^{2k}/(4^k(2k+1)!)$,
\[
-w\big(S(w\hbar\partial_p)-1\big)\Wv(p)=\sum_{k\geq 1}w^{2k+1}\hbar^{2k}\epsilon_k(p),\qquad \epsilon_k:=-\frac{\Wv^{(2k)}}{4^k(2k+1)!}\quad(\text{local letters}),
\]
with $\Wv^{(2k)}=(p^4/4)^{(2k)}+(2k)!\,v\,p^{-2k-1}\in\bbQ[v][p,p^{-1}]$. In $W_i$ every $\hbar^{2k}$ comes with $w_i^{2k+1}$, in $W_{i,j}$ every $\hbar^2$ with $(w_i+w_j)^2$ and a factor $(p_i-p_j)^{-2}$; so at fixed order in $\hbar$ the degrees in the $w_i$ are bounded, finitely many $r_i$ and graphs contribute, and the coefficients are polynomials in the $\epsilon_k(p_i)$ and the $(p_i-p_j)^{-1}$. Finally $\mathcal D_x(f\,\dd p)=\dd\big(p^2f/(p^5-v)\big)$ only adds powers of $(p^5-v)^{-1}$. (2) \emph{Clearing.} Write $\omega_{h,n}=P\prod_i\dd p_i\big/\big(\prod_ip_i^{a}(p_i^5-v)^{K}\prod_{i<j}(p_i-p_j)^{e}\big)$ with $P\in\bbQ[v][p]$. By \cref{conv:mod-kernel} there are no poles at $p_i=0$ or on the diagonals, so for every $v\neq 0$, $P$ vanishes to order $\geq e$ along $p_i=p_j$ and $\geq a$ along $p_i=0$. Dividing by the monic polynomials $(p_i-p_j)^e$, $p_i^a$ in $p_i$ over $\bbQ[v][p_k:k\neq i]$, the remainder vanishes for all $v\neq 0$ and $p$, hence is $0$; one factor at a time, these denominators cancel in $\bbQ[v][p]$, without negative powers of $v$.
\end{proof}

\begin{thm}[L-regularity]
\label{thm:mod-Lreg}
Let $2h-2+n>0$. For legs $\phi_1,\dots,\phi_n\in\bbQ[v][p]$ the residue pairing $\langle\phi_1\cdots\phi_n\rangle^{\mathrm{TR}}_h$ lies in $\bbQ[v]$. For legs in $\cL=\bbQ[v][p]+\bbQ(v)[x]$ the residue pairing lies in $\bbQ[v]$; the components in $\bbQ(v)[x]$ pair to zero.
\end{thm}

\begin{proof}
By \cref{lem:mod-abdks}, $\phi_l(p_l)\omega_{h,n}$ is a rational $1$-form in $p_l$ with poles at the branch points and possibly at $\infty$, holomorphic at $p_l=0$; so $\sum_b\Res_{p_l=\varsigma_b}=-\Res_{p_l=\infty}$. Expanding $(p^5-v)^{-K}=\sum_{j\geq 0}\binom{K+j-1}{j}v^jp^{-5K-5j}$ for $|p|^5>|v|$, $\Res_{p=\infty}p^a\,\dd p/(p^5-v)^K$ is $-\binom{K+j-1}{j}v^j$ if $a+1=5(K+j)$ with an integer $j\geq 0$, and $0$ otherwise, and the result is again of the form of \cref{lem:mod-abdks} in the other variables. The iterated residues commute (no diagonal poles), so
\[
\langle\phi_1\cdots\phi_n\rangle^{\mathrm{TR}}_h=(-1)^n\Res_{p_1=\infty}\cdots\Res_{p_n=\infty}\Big[\prod_l\phi_l(p_l)\,N_{h,n}\prod_l\dd p_l\Big/\prod_l(p_l^5-v)^K\Big]\in\bbQ[v].
\] The components in $\bbQ(v)[x]$ pair to zero by \cref{lem:mod-xfun}.
\end{proof}

\begin{prop}[The period of $\omega_{h,1}$]
\label{prop:mod-period}
Let $h\geq 1$. Then:
\begin{enumerate}
\item $\omega_{h,1}=\dd\Phi_h$ with $\Phi_h$ rational and holomorphic at $p=0$ and $p=\infty$;
\item $\int_\infty^0\omega_{h,1}=\frac{B_{2h}}{2h}\,v^{1-2h}$ along any path avoiding the branch points;
\item for $h\geq 2$, $\FTR_h=\dfrac{B_{2h}}{2h(2h-2)}\,v^{2-2h}$.
\end{enumerate}
\end{prop}

\begin{proof}
\emph{Step 1 (exactness).} Put $\omega^{\mathrm A}_{h,1}:=\omega^{(h)}_{1,0}=-\omega_{h,1}$. By \eqref{eq:mod-710}, $\omega^{\mathrm A}_{h,1}=[\hbar^{2h}]\sum_{r\geq 1}\mathcal D_x^r[w^r]W_1$; the term $r=0$ is $-[w^1]e^{E}\dd p=0$, as $E:=\sum_kw^{2k+1}\hbar^{2k}\epsilon_k=O(w^3)$. For $r\geq 1$, $\mathcal D_x^r(a\,\dd p)=\dd\big(D_x^{r-1}(a/x')\big)$ (\cref{def:mod-xi}), so $\omega^{\mathrm A}_{h,1}=\dd\Phi^{\mathrm A}_h$ with $\Phi^{\mathrm A}_h:=\sum_{r\geq 1}D_x^{r-1}(a_r/x')$, $a_r:=-[\hbar^{2h}][w^{r+1}]e^{E}=-\sum\frac1{j!}\prod_{i=1}^j\epsilon_{k_i}$, the finite sum over $1\leq j\leq h$ and $k_1+\dots+k_j=h$ with $r=2h-1+j$. Put $\Phi_h:=-\Phi^{\mathrm A}_h$.

\emph{Step 2 ($p=\infty$).} $\deg\epsilon_k\leq 4-2k$, $\deg(1/x')=-3$, and $D_x$ lowers the degree by $4$, so every term of $\Phi^{\mathrm A}_h$ has degree at most $(4j-2h)-3-4(2h-2+j)=5-10h<0$: $\Phi^{\mathrm A}_h(\infty)=0$.

\emph{Step 3 ($p=0$).} The \emph{pure part} of a term replaces $\epsilon_k$, $1/x'=-(p^2/v)(1-p^5/v)^{-1}$ and $D_x$ by their leading parts $\epsilon^0_k:=-v\,p^{-2k-1}/(4^k(2k+1))$, $-p^2/v$ and $\delta_0:=-(p^2/v)\,\dd/\dd p$; in $\bar p=v/p$, $\delta_0=\dd/\dd\bar p$, $-p^2/v=-v/\bar p^2$, $\epsilon^0_k=-\bar p^{2k+1}v^{-2k}/(4^k(2k+1))$. With the local letters $\tilde w:=w\bar p$, $\hbar':=\hbar/(2v)$, $E_0:=-\sum_{k\geq 1}\hbar'^{2k}\tilde w^{2k+1}/(2k+1)$ and $C_{h,r}:=[\hbar'^{2h}][\tilde w^{r+1}]\exp(E_0)$, the pure part of $a_r$ is $-(2v)^{-2h}\bar p^{r+1}C_{h,r}$, and that of the $r$-th term is $(\dd/\dd\bar p)^{r-1}\big[(2v)^{-2h}v\,C_{h,r}\bar p^{r-1}\big]=v(2v)^{-2h}(r-1)!\,C_{h,r}$, a constant. Replacing a leading part by the full factor adds terms of order higher in $p$ by at least $3$ (by $\geq 2k+1$ for $\epsilon_k$, $\geq 5$ for $1/x'$ and for $D_x$), and $\delta_0$ never lowers it; so $\Phi^{\mathrm A}_h$ is holomorphic at $0$ with $\Phi^{\mathrm A}_h(0)=v(2v)^{-2h}\Sigma_h$, $\Sigma_h:=\sum_r(r-1)!\,C_{h,r}$.

\emph{Step 4 ($\Sigma_h=-4^hB_{2h}/(2h)$).} Let $\mathsf a(\mathsf x):=\sum_{k\geq 1}\mathsf x^k/(2k+1)$, so $E_0=-\tilde w\,\mathsf a(\hbar'^2\tilde w^2)$ and, with $r+1=2h+j$, $C_{h,r}=\frac{(-1)^j}{j!}[\mathsf x^h]\mathsf a(\mathsf x)^j$. By the binomial series (the terms $j=0$ and $j>h$ have no $\mathsf x^h$),
\[
\Sigma_h=\sum_{j=1}^h(-1)^j\frac{(2h-2+j)!}{j!}[\mathsf x^h]\mathsf a^j=(2h-2)!\,[\mathsf x^h](1+\mathsf a(\mathsf x))^{1-2h}.
\]
With $\mathsf x=\mathsf y^2$, $1+\mathsf a(\mathsf y^2)=\operatorname{artanh}(\mathsf y)/\mathsf y$, so $[\mathsf x^h](1+\mathsf a)^{1-2h}=[\mathsf y^{1}]\operatorname{artanh}(\mathsf y)^{1-2h}$. The Lagrange--B\"urmann formula for $\mathsf y=\tanh\mathsf u$, $[\mathsf y^n]H(\mathsf u(\mathsf y))=\frac1n\Res_{\mathsf u=0}H'(\mathsf u)\tanh(\mathsf u)^{-n}\dd\mathsf u$ ($n\neq 0$, $H$ a Laurent series; change of variables and one integration by parts), with $H=\mathsf u^{1-2h}$, $n=1$, gives $(1-2h)[\mathsf u^{2h}](\mathsf u\coth\mathsf u)=(1-2h)4^hB_{2h}/(2h)!$. So $\Sigma_h=-4^hB_{2h}/(2h)$ and $\Phi^{\mathrm A}_h(0)=-\frac{B_{2h}}{2h}v^{1-2h}$. This proves (i).

\emph{Step 5.} (ii) $\int_\infty^0\omega_{h,1}=\Phi_h(0)-\Phi_h(\infty)=-\Phi^{\mathrm A}_h(0)$. (iii) As $\Phi_{y\dd x}=p^5/5-v\log p$, $(2-2h)\FTR_h=\frac15\langle p^5\rangle^{\mathrm{TR}}_h+v\langle-\log p\rangle^{\mathrm{TR}}_h$ (\cref{def:mod-pair}). Here $\langle p^5\rangle^{\mathrm{TR}}_h\in\bbQ[v]$ (\cref{thm:mod-Lreg}) has exponent $2-2h<0$ (\cref{lem:mod-scale}~(ii)), hence vanishes, and $\langle-\log p\rangle^{\mathrm{TR}}_h=\int_{\mathfrak r}\omega_{h,1}=-\frac{B_{2h}}{2h}v^{1-2h}$ by (ii). So $\FTR_h=\frac{-v}{2-2h}\cdot\frac{B_{2h}}{2h}v^{1-2h}$.
\end{proof}

\begin{thm}[The Gap Theorem]
\label{thm:mod-gap}
On the model curve, for every $h\geq 2$, $\FTR_h=\dfrac{B_{2h}}{2h(2h-2)}v^{2-2h}$ in the convention of \cref{conv:mod-kernel}; hence $\Fm_h=c_hv^{2-2h}$ with $c_h=(-5)^{1-h}\dfrac{B_{2h}}{2h(2h-2)}=\dfrac{(-1)^{h-1}B_{2h}}{2h(2h-2)\,5^{h-1}}$.
\end{thm}

\begin{proof}
\Cref{prop:mod-period}~(iii) and \cref{cor:mod-trF}.
\end{proof}

\begin{rmk}[A first proof through a formula of Hock]
\label{rmk:mod-hock}
The formula \cite[Cor.~1.2, (1.3)]{Hock25}, whose conventions \cite[(2.1), (1.1)]{Hock25} are those of \cref{conv:mod-kernel}, also gives $\FTR_h$: for $y=p$ only the leading terms at $p=0$ contribute, and Step~4 above applies. This route is not used. \draftnote{The closed forms printed in \cite[\S3.2]{Hock25} are not used; only Cor.~1.2 there is cited.}
\end{rmk}

\input{sections/s06_model_8}



